\section{经验}

在构建和部署GFS的过程中，我们遇到过各种问题，其中一些属于运维问题，另一些属于技术问题。

起初，GFS被设想为生产系统的后端文件系统。随着时间推移，其使用范围扩展到研究与开发任务。
它最初几乎不支持权限、配额等功能，但现在已经具备这些功能的初步形式。生产系统通常纪律
严明、管理受控，但用户有时并非如此。我们需要更多基础设施来防止用户彼此干扰。

我们遇到的一些最大问题与磁盘和Linux有关。许多磁盘向Linux驱动程序声明支持一系列
IDE协议版本，但实际上只有在较新的版本下才能可靠响应。由于这些协议版本非常相似，
这些硬盘大多数时候可以正常工作；但偶尔，版本不匹配会导致硬盘与内核对硬盘状态产生分歧。
内核中的问题会因此悄无声息地损坏数据。这个问题促使我们使用校验和来检测数据损坏；
与此同时，我们修改了内核，以处理这些协议不匹配的情况。

此前，由于\textbf{fsync()}的开销，我们在Linux2.2内核上遇到过一些问题。其开销与文
件大小成正比，而不是与被修改部分的大小成正比。对于我们的大型操作日志而言，这
尤其是在实现checkpoint之前是一个问题。我们一度通过使用同步写入来规避它，最终迁
移到了Linux2.4。

另一个Linux问题是单个读写锁：地址空间中的任何线程从磁盘换入页面时需要持有它(读锁)，
或者在\textbf{mmap()}调用中修改地址空间时需要持有它(写锁)。我们曾在轻负载下观察到系统出
现瞬态超时，并仔细排查资源瓶颈或偶发的硬件故障。最终，我们发现，当磁盘线程正在换
入先前已映射的数据时，这把单独的锁会阻塞主网络线程将新数据映射到内存中。由于我们
的主要限制来自网络接口，而非内存拷贝带宽，我们以额外一次拷贝为代价，
将\textbf{mmap()}替换为\textbf{pread()}，从而规避了这个问题。

尽管偶尔会遇到问题，Linux代码的可用性仍多次帮助我们探索并理解系统行为。在适当情
况下，我们会改进内核，并将这些改动分享给开源社区。
